%% ============================================================================
%% sections/02-background.tex -- Background and Motivation
%% Class-independent section file: it is \input by main-acm.tex and main-ieee.tex.
%% Rules (see README-build.md): no class-specific commands; cite with \cite only;
%% figure sizes come from the float budget and table bodies from tables/, never from this file.
%% Every paragraph and every caption is one numbered unit that starts with a
%% marker line "%% P:<id>".  Ids are frozen: never renumber, never reuse.
%% The same ids appear in paper.zh-CN.md; check_align.py compares the two.
%% Word budget of this section: 690 words of running text.
%% ============================================================================

\section{Background and Motivation}\label{sec:background}

\subsection{Block-Scaled Formats and Three Block Axes}\label{sec:bg-formats}

%% P:sec2.1-p1
\ptag{sec2.1-p1}%
A \blockscaled format lets $k$ elements $P_i$ share one scale $X$; element $i$ represents $X \cdot P_i$ \cite{OCP2023MXSpec}.
We use the term for three families: block floating point (BFP), where integer mantissas share an exponent and the block size varies by design \cite{Fowers2018Brainwave,Drumond2018HBFP,DarvishRouhani2020MSFP}; OCP MX, with $k=32$ and an E8M0 scale \cite{OCP2023MXSpec,DarvishRouhani2023MXFormats}; and NVFP4, with E2M1 elements, an E4M3 scale per 16 elements and an FP32 scale per tensor \cite{NVIDIA2025NVFP4Blog}.

%% P:sec2.1-p2
\ptag{sec2.1-p2}%
We keep three block axes apart; values are the prototype's.
The input scale block runs along the inner-product dimension $K$: in the MAC cells' MX input mode, 32 elements share an E8M0 scale.
The output block runs along the output dimension $N$: a row's 16 elements share one power-of-two exponent.
The storage group is a row of the \sidecar, which stores block exponents apart from the payload, 16 per row.
The GEMM outlet thus forms a new block representation and passes no input scale through \TBD{D7}{how the three axes read in the target design}.

\subsection{Where the Hand-off Happens Today}\label{sec:bg-today}

%% P:sec2.2-p1
\ptag{sec2.2-p1}%
We survey the designs we examined whose public material suffices to determine the \handoff format.
On NVIDIA Blackwell, \blockscaled products accumulate in FP32 and cuBLAS quantizes at the end of the GEMM \cite{NVIDIA2026PTXISA94,nvidia_cublas_13_4}.
MTIA~2i forwards FP32 results directly to a floating-point SIMD engine \cite{Coburn2025MTIA2}.
Tenstorrent stores BFP; its default element-wise path unpacks to FP32 or BF16 before compute and packs on write-back \cite{tenstorrent_isa_wormhole,tenstorrent_metalium_dataflow}.
Trainium3 and Maia~200 output scalars; block formats are rebuilt downstream \citeTBD{Trainium3; Maia 200}.
Outside matrix multiplication, Intel AMX, PTX and CDNA~4 accept a block scale only in format-conversion instructions \cite{NVIDIA2026PTXISA94,amd2025cdna4} \citeTBD{Intel AMX}.

%% P:sec2.2-p2
\ptag{sec2.2-p2}%
We do not claim that these designs could not keep a block representation, and we do not dismiss prior work as special-purpose; we ask under which conditions the placement of the \handoff changes cost and numerics.
Two surveyed practices are strong baselines: the fused end-of-GEMM quantization is the production counterpart of \armSE, and direct matrix-to-vector forwarding is established practice, so \armBD differs only in the forwarded format.
What public material does not disclose is not assumed absent.

\subsection{Two Hypotheses to Test}\label{sec:bg-hypotheses}

%% P:sec2.3-p1
\ptag{sec2.3-p1}%
We pose two hypotheses, not conclusions.
Hypothesis~A: fusion placement, result lifetime and the point of numeric narrowing together decide whether a block-form \handoff pays off. For a local, single-consumer operation, producer-side fusion wins under some resource conditions, whereas a result that is retained, consumed repeatedly, or held for cross-block information may favor the block form in a predictable region.
Hypothesis~B: metadata transaction granularity can decide whether keeping a block form resident helps or hurts, with the numerical representation unchanged.
We test them separately and combine them in \secref{sec:res-payoff}, with narrowing points fixed per placement (\secref{sec:placements}).

\subsection{A Workload in Motion and the Design Questions It Raises}\label{sec:bg-workload}

%% P:fig2-cap
\FigSlot{fig:workload}%
  {Share of multiply--accumulate operations by model component against context length, for the three models}%
  {source plot exists; to be re-laid out for one column}%
  {Share of multiply--accumulate operations per decoded token, by model component, versus context length (\evAnalytic). V4.1-Flash at 64K and above follows its report's count.}

%% P:tab1-cap
\begin{table}[t]
\centering\SlotName{tab:workload-gen}%
\caption{Demands of one decode step for three models from two generations; analytical counts, not measurements. V4.1-Flash at 1M follows its report's count; cache bytes include block scales.}\label{tab:workload-gen}
%% tables/tab-workload-gen.tex -- body of the table `tab:workload-gen'.
%% The float, the caption and the label are in sections/02-background.tex; this file sets the type size and the
%% column separation itself, so that it alone decides how wide and how high the table is.
%% No blank line in this file: check_align.py reads it as part of the unit that inputs it.
\TableSize\setlength{\tabcolsep}{3pt}%
\begin{tabular}{@{}lrrr@{}}
\toprule
Per decoded token & V4-Pro & V4-Flash & V4.1-Flash \\
\midrule
Layers / hidden width & 61 / 7,168 & 43 / 4,096 & 40 / 5,120 \\
MACs at 8K context ($10^{10}$) & 5.466 & 1.477 & 1.796 \\
MACs at 1M context ($10^{10}$) & 15.16 & 7.017 & 2.875 \\
Largest at 8K: feed-forward (\%) & 51.6 & 51.3 & 55.2 \\
Largest at 1M: index scoring (\%) & 42.5 & 64.3 & 38.3 \\
Cache bytes read at 1M (MB) & 705.3 & 479.3 & 183.3 \\
Main key--value entry, 512 dims & \multicolumn{2}{c}{448 8-bit + 64 BF16} & 512 4-bit \\
8-bit blocks: weights / activations & \multicolumn{2}{c}{128$\times$128 / 128} & 32$\times$32 / 32 \\
\midrule
\multicolumn{4}{@{}l@{}}{All three: one ReLU and weighted add per 128 index-scoring MACs;} \\
\multicolumn{4}{@{}l@{}}{one softmax-accumulator rescale per 128 attention-core MACs.} \\
\bottomrule
\end{tabular}

\end{table}

%% P:sec2.4-p1
\ptag{sec2.4-p1}%
The three models share a context limit of one million tokens.
We derive the multiply--accumulate operations (MACs) per decoded token analytically from their public configurations and reference implementations (\secref{sec:workload-method}), and group them by model component: index scoring (an indexer's 128-dimensional inner products with cached keys, which select the top-$k$ entries), the attention core (scores and weighted sum of the attended entries under an online softmax) and the feed-forward experts.
The mix shifts with context (\figref{fig:workload}, \tabref{tab:workload-gen}): at 8K tokens the feed-forward experts hold 51.3--55.2\% of the MACs; at 1M index scoring is the largest component, at 38.3--64.3\% (V4.1-Flash by its technical report's count, \secref{sec:workload-method}).

%% P:sec2.4-p2
\ptag{sec2.4-p2}%
Structural quantities held across the three models: head dimension 512, index dimension 128 and an online softmax in groups of 64 entries.
Operation shares, storage formats and scale-block sizes changed: V4.1-Flash stores its main key--value entries in 4-bit form and shrinks 8-bit blocks from 128 to 32.
Its authors state that this format saves storage rather than speeding up matrix multiplication \cite{DeepSeekAI2026V41Flash}; the V4 authors propose replacing SwiGLU with an element-wise activation free of exponentials and divisions, so that post-GEMM processing does not stall the GEMM pipeline \citeTBD{DeepSeek-V4 paper}.

%% P:sec2.4-p3
\ptag{sec2.4-p3}%
Six design questions follow; the target design's answers are planning points, and \secref{sec:uarch} states what the prototype has.
(i) Which unit executes non-GEMM operations: two \VectorCore{}s per PE, their form open (\secref{sec:consumer}).
(ii) One PE type or two: one, with capability bits for a second.
(iii) Compute per bit-width tier: 2048, 1024 and 256 MACs per PE per cycle at 4-bit, 8-bit and BF16.
(iv) What a \blockscaled \handoff saves: this paper's question.
(v) Shared storage and PE port: 128--256\,MiB per $8\times8$ mesh and a local port of 64 payload bytes per cycle into each PE (\secref{sec:pe-org}).
(vi) Fixed cost per tile, one pass of the array: 5 cycles.
